% !TEX root = ../notes_missingsemester.tex
% Block 9: Watch It Run — Monitoring, Logging & Analytics
%%%%%%%%%%%%%%%%%%%%%%%%%%%%%%%%%%%%%%%%%%%%%%%%%%%%%%%%%%%%%%%%%%%%%%

\index{monitoring}\index{logging}\index{analytics}\index{Discord webhook}
\chapter{Watch It Run: Monitoring, Logging \& Analytics}
\printchapterversion{09_monitoring_logging}
\newthought{You cannot improve what you cannot measure.}

% ==============================================================================================
% BLOCK 1: INTRODUCTION
% ==============================================================================================

\section{From Debugging to Observability}

\newthought{You already know \texttt{journalctl}} for reading logs and debugging crashes (Block~5).
Now the question changes: your service is running and nobody is watching the terminal. How do you know
it is working well? Three pillars:

\begin{description}
    \item[Logs] Structured, queryable records of what happened.
    \item[Monitoring] Automated checks that alert you when something breaks.
    \item[Analytics] Patterns in your data that reveal how the system behaves over time.
\end{description}

\newpage
\section{Optional Exercise: Catch Up to \texttt{v0.7}}\index{exercises}\index{Git!tag}\index{recovery}

\newthought{Catching up via snapshot and tag.}\index{Git!tag}\index{recovery}
This block builds on the state at the end of Block~8: GitHub Actions
workflows under \texttt{.github/workflows/} that lint and test on every
push and deploy to \texttt{prod} on every merge to \texttt{main},
GitHub repository secrets configured for the deploy SSH key, a deploy
script in \texttt{scripts/deploy.sh}, cron jobs installed on \texttt{prod}
for nightly backups and weekly \texttt{pip-audit} runs, and feature flags
read from environment variables in \texttt{app/main.py}.
If you joined late, missed Block~8, or your repository has drifted,
restore from the snapshot and reset to \texttt{v0.7} rather than redoing
every prior exercise.

\marginnote{\newthought{Tag map.}\index{Git!tag}
\texttt{v0.7} marks the end of Block~8, \texttt{v0.8} will mark the end
of this block.
Browse the full set at
\url{https://github.com/schutera/pixelwise/tags}.}

\newthought{Restore the snapshot.}
On both VMs, restore the \texttt{B8-complete} snapshot from Block~8.
The snapshot carries everything that is not in Git, the \texttt{.venv/},
the model artefact, the systemd service, the PostgreSQL database, the
Nginx configuration, the TLS certificate, the cron table on \texttt{prod},
and any real \texttt{.env} values you wrote.
GitHub-side state, the repository secrets and the workflow runs, lives
on GitHub itself and is not part of the VM snapshot.
With the system layer back in place, only the repository state needs
to be rewound.

\newthought{Catch-up procedure on \texttt{dev}.}
Reset PixelWise to the previous tag:

\begin{verbatim}
cd ~/pixelwise
git fetch origin
git reset --hard v0.7
\end{verbatim}

\marginnote{\newthought{Why \texttt{reset -{}-hard}?}\index{Git!reset}
This rewrites your working tree to match the end of Block~8 exactly,
including \texttt{.github/workflows/}, \texttt{scripts/deploy.sh}, the
backup scripts, and the feature-flag wiring in \texttt{app/main.py}.
There is no work to lose if you have not committed beyond \texttt{v0.7};
if you have, branch off first with \texttt{git branch backup-pre-catchup}
so the work is recoverable.}

\newthought{Mirror it on \texttt{prod}.}

\begin{verbatim}
ssh produser@192.168.56.11
cd /opt/pixelwise
git fetch origin
git reset --hard v0.7
sudo systemctl reload nginx
sudo systemctl restart pixelwise
\end{verbatim}

\marginnote{\newthought{Already have a fork?}\index{Git!fork}
If your own fork carries the \texttt{v0.7} tag, swap \texttt{origin} to
your fork URL on \texttt{dev} with
\texttt{git remote set-url origin git@github.com:<you>/pixelwise.git}
so push and pull continue to target your account.
The HTTPS clone on \texttt{prod} stays read-only.
Re-create the GitHub repository secrets on the fork as well, since they
do not transfer.}

\newthought{Verify the automation.}
Trigger a no-op commit and watch the CI workflow run on GitHub:

\begin{verbatim}
git commit --allow-empty -m "Verify CI catch-up"
git push
\end{verbatim}

The Actions tab should show a green run for the lint and test workflow.
On \texttt{prod}, confirm the cron jobs are still installed:

\begin{verbatim}
crontab -l
\end{verbatim}

You are now at the state where Block~8 ended, with both machines and
GitHub aligned, ready to start Block~9.

\newthought{Without a snapshot,} replay Block~8 before continuing.
The file state alone is not enough when the next exercises expect
working CI/CD, scheduled cron jobs on \texttt{prod}, and feature flags
that the runtime can already read.

% ==============================================================================================
% BLOCK 2: CONCEPTS AND PRACTICE
% ==============================================================================================

\newpage
\section{Structured Logging}\index{structured logging}\index{JSON logging}

\newthought{Plain strings are not enough.} \texttt{print("got request")} tells you something happened.
It does not tell you what, when, from whom, how long it took, or what the result was. You cannot filter,
aggregate, or search plain strings at scale.

\newthought{Structured logging} means every log line is a JSON dict:

\begin{verbatim}
{
  "timestamp": "2026-04-05T14:23:01Z",
  "level": "INFO",
  "endpoint": "/classify",
  "request_id": "abc-123",
  "latency_ms": 42,
  "predicted_class": "7",
  "confidence": 0.94,
  "model_version": "v1"
}
\end{verbatim}

\marginnote{Structured logs are parseable by machines, readable by humans, and queryable with \texttt{jq}\index{jq}
(\url{https://jqlang.github.io/jq/})
or piped into any log aggregation tool.\\
Python \texttt{logging}\index{logging}: \url{https://docs.python.org/3/library/logging.html}}

\newthought{In Python,} the standard \texttt{logging} module accepts a custom formatter that overrides
\texttt{format(record)} to serialise each record as JSON. The formatter pulls the timestamp, level,
and message from the record itself, then copies known extra fields if the calling code attached them
through \texttt{logger.info(..., extra=\{...\})}. The exercise below wires the formatter into
\texttt{app/main.py} and lists the fields PixelWise emits.

\newpage
\section{Exercise: Add Structured JSON Logging}\index{exercises}

\marginnote{If your VM is broken, restore the \texttt{B8-complete} snapshot.
At the end of this session, take a snapshot named \texttt{B9-complete}.}

\newthought{Configure the Python \texttt{logging} module} with a JSON formatter using the class shown
in the theory above. Every request should log: \texttt{timestamp}, \texttt{level}, \texttt{endpoint},
\texttt{request\_id}, \texttt{latency\_ms}, \texttt{predicted\_class}, \texttt{confidence},
\texttt{model\_version}.

\newpage
\section{The \texttt{/health} Endpoint Upgraded}\index{health check}

\newthought{Block~5's \texttt{/health}} confirmed the process is running. Now distinguish two levels:

\begin{description}
    \item[Liveness] Is the process running? Returns 200 if yes.
    \item[Readiness] Can it serve requests? Is the DB reachable? Is the model loaded?
\end{description}

The upgraded handler runs both checks in sequence, opens a short test query against PostgreSQL,
and confirms that the global model pipeline is loaded. If either fails, the response status flips
from \texttt{"ok"} to \texttt{"degraded"} so the cron-driven monitor can react. The exercise below
has the full handler PixelWise needs.

\marginnote{What to expose externally vs what to keep internal: the public \texttt{/health} endpoint
returns \texttt{"ok"} or \texttt{"degraded"}. It does not expose DB passwords, stack traces, or version
details. Detailed diagnostics are logged internally.}

\newpage
\section{Exercise: Upgrade \texttt{/health}}\index{exercises}

\newthought{Check DB connectivity and model load status.} Return \texttt{\{"status": "ok"\}} externally;
log detailed diagnostics internally.

\newpage
\section{Monitoring with Discord Webhooks}\index{Discord webhook}\index{monitoring}

\newthought{No new tools to install.} A cron job \texttt{curl}s \texttt{/health} every 5 minutes.
If the response is not \texttt{"ok"}, it sends a Discord webhook notification. Your service talks to you
when something breaks, on your phone, in your pocket.

\newthought{The shape of the script.} A short bash file under \texttt{scripts/health\_check.sh}
calls \texttt{curl} against the local \texttt{/health} endpoint, parses the \texttt{status} field
with \texttt{jq}, and posts a JSON payload to the webhook URL when the status is anything other
than \texttt{"ok"}. The webhook URL itself comes from the environment, so the script carries no
secrets in the repository. The exercise below has the full script and the cron registration.

\marginnote{The Discord webhook URL is stored in \texttt{.env}. It is a secret that grants posting
access to your channel. The cron job runs the script; the script talks to Discord when there is a problem.\\
Discord Webhooks: \url{https://discord.com/developers/docs/resources/webhook}}

\newpage
\section{Exercise: Set Up Discord Webhook Monitoring}\index{exercises}

\newthought{Wire the alerting path end to end.} Work through the steps below:

\begin{itemize}
    \item Create a Discord server and webhook (Settings $\to$ Integrations $\to$ Webhooks).
    \item Add \texttt{DISCORD\_WEBHOOK\_URL} to \texttt{.env} and \texttt{.env.example}.
    \item Write \texttt{scripts/health\_check.sh} using the script from the theory section.
    \item Register the cron job:
    \begin{verbatim}
    crontab -e
    # Add:
    */5 * * * * /opt/pixelwise/scripts/health_check.sh
    \end{verbatim}
    \item Test it: stop the \texttt{systemd} service and wait. Does Discord ping?
    \begin{verbatim}
    sudo systemctl stop pixelwise
    # Wait up to 5 minutes, check Discord
    sudo systemctl start pixelwise
    \end{verbatim}
\end{itemize}

\newpage
\section{Event-Driven Notifications}\index{alerting}

\newthought{The Discord pattern generalises.} Any event worth knowing about can trigger a notification.
A second example sits inside the classify endpoint: After the model returns a result, compare its
confidence against a threshold constant, and call a small \texttt{send\_discord\_alert} helper with the
predicted class and the confidence score whenever the value falls short. The exercise below wires it up.

\newthought{If these spike,} the model may be struggling with inputs it was not trained on, a real signal,
not noise. This is the natural step toward model monitoring: tracking confidence over time, spotting
distribution shifts, detecting when a model goes stale.

\newpage
\section{Exercise: Add a Low-Confidence Alert}\index{exercises}

\newthought{In the classify endpoint,} send a Discord notification when confidence falls below a
threshold (e.g., 0.5). Test it by drawing an ambiguous digit. Does Discord ping?

\newpage
\section{Exercise: Add a Canvas-Clear Notification}\index{exercises}

\newthought{Add a second Discord notification} that fires when a user clears the drawing field.
Add an API endpoint (e.g., \texttt{POST /api/events/clear}) that the frontend calls on canvas clear.
Discuss: what does a spike in clears tell you about model performance?

\newpage
\section{Analytics}\index{analytics}

\newthought{Everything that happens in the app is a potential data point:}

\begin{description}
    \item[Backend signals] Prediction counts, confidence values, model version usage, already stored in
    PostgreSQL from Block~6.
    \item[Frontend signals] Page loads, drawings submitted, canvas clears, time spent drawing. Any user
    interaction can be tracked by the frontend and sent to the API.
\end{description}

\subsection{Analytics Concepts}

\marginnote{These concepts are how product teams understand user behaviour at scale:
\textbf{funnels}\index{funnel} (open page $\to$ draw $\to$ submit $\to$ view result, where do users drop off?),
\textbf{retention} (do users come back?),
\textbf{sessions}\index{session} (grouping events into a single visit, recall the session IDs from Block~7's URL design).}

\subsection{The \texttt{/analytics} Endpoint}\index{/analytics}

\newthought{Provided as boilerplate.} The endpoint queries the \texttt{predictions} table and returns
total prediction counts, a per-class breakdown, the average confidence value, and a tally of which model
versions served traffic. Students hit the endpoint and see the data. The lesson is not how to write the
queries, but what the numbers tell you about your app. The exercise below has the full handler and the
\texttt{curl} call that exercises it.

\marginnote{In production, analytics panels are never public. They reveal usage patterns and internal metrics.
They belong behind authentication on a separate internal route, for example \texttt{/admin/analytics}.}

\subsection{Frontend Analytics}

\newthought{Instrument the frontend} to count user interactions. A small panel in the UI displays:
page loads, drawings submitted, canvas clears. These are frontend-only counters, simple JavaScript
variables incremented on each event.

\newthought{Canvas clears are a signal too.} When a user clears the drawing field, it often means the
prediction was wrong or the user wants to try again. This is the natural step to wire in model analytics.
A spike in clears correlated with low-confidence predictions tells you the model is struggling with certain inputs.

\newpage
\section{Exercise: Deploy the \texttt{/analytics} Endpoint}\index{exercises}

\newthought{Use the provided boilerplate.} Hit it with \texttt{curl} and interpret the results:

\begin{verbatim}
    curl https://192.168.56.11/api/analytics -k \
      -H "X-API-Key: your-key"
\end{verbatim}

\newthought{Read the numbers.} Which digit is drawn most? Is one model version more confident than the other?

\newpage
\section{Exercise: Add Frontend Analytics}\index{exercises}

\newthought{Instrument the starter code} to count page loads, drawings submitted, and canvas clears.
Display these in a small panel in the UI. Discuss: what would you track if this were a real product?

\newpage
\section{Exercise: Commit}\index{exercises}

\newthought{Stage and push} the new logging, alerting, and analytics work:

\begin{verbatim}
    git add app/ scripts/ frontend/ .env.example
    git commit -m "Add structured logging, Discord alerts, analytics"
    git push
\end{verbatim}

\newpage
\section{Exercise: Take a Snapshot}\index{exercises}

\newthought{Snapshot both VMs} and name them \texttt{B9-complete}.

\newthought{A note on what comes next.} With the service observable end to end, the focus shifts from
watching the system to scaling and hardening it, a thread we will revisit later in the course.

\newpage
\section{Security Thread}

\newthought{Log sanitization.} Never log raw pixel data or user identifiers unnecessarily.
The \texttt{/health} endpoint confirms the system is alive but does not expose internal state
(DB passwords, stack traces, version details) to unauthenticated callers.

% ==============================================================================================
% BLOCK 3: CONSOLIDATION
% ==============================================================================================

\section{Self-Reflection and Recap}

\newthought{Reflection questions:}
\begin{itemize}
    \item What is the difference between plain-text logging and structured logging? Why does it matter?
    \item Why do we monitor with Discord webhooks instead of installing a monitoring platform?
    \item What does a spike in low-confidence predictions tell you about your model?
    \item Why should analytics endpoints never be public?
    \item What would a spike in canvas clears correlated with low-confidence predictions indicate?
\end{itemize}

\newthought{Milestone:} PixelWise is fully observable. Logs are structured and queryable. Discord alerts
on health check failure and low-confidence predictions. Analytics endpoint and frontend counters show
usage patterns. Students can answer: ``how is my app being used, and how is my model performing?''
